\input{./._relpath-to-latexroot.ltx}
\documentclass[\latexroot/\projectname]{subfiles}
\input{\latexroot/@resources/subfile-setup.ltx}

\makeatletter
\renewenvironment{verbatimwrite}[1]{\comment}{\endcomment}
\makeatother

\begin{document}

\whenintegrated{%
  \input{\latexroot/._relpath-to-latexroot.ltx}%
}

\section{The Model}\whenintegrated{\label{model}}
\whenintegrated{\label{sec:model}}

This section explains the theoretical model objects that I reproduce in the notebook. The aim is not to rewrite the original Cagetti and De Nardi paper, but to identify the pieces of the household problem that have to be represented in code: the state variables, the worker branch, the entrepreneur branch, the enforcement constraint, and the continuation values.

\subsection{Household State and Timing}

There is a continuum of households and no aggregate uncertainty. The full individual state can be written as
\[
  \mathbf{x} = (a,y,\theta,s),
\]
where $a$ is beginning-of-period assets, $y$ is worker productivity, $\theta$ is entrepreneurial ability, and $s$ records the household's demographic and occupational status. In the paper this status separates young workers, young entrepreneurs, old entrepreneurs, and old retirees. In the notebook I keep the same economic distinction through the value functions $V$, $V_w$, $V_e$, $W$, $W_e$, and $W_r$.

The timing is useful to state explicitly. At the start of the period, the household observes assets and current abilities. It then makes a discrete status decision: a young household chooses between working and entrepreneurship, while an old entrepreneur chooses between continuing the business and retiring. After that branch is chosen, the household chooses continuous controls. Everyone chooses consumption $c$ and next-period assets $a'$. Entrepreneurs also choose working capital $k$, which is an intra-period business scale and does not itself carry into the next period as a state variable.

\subsection{Preferences, Demographics, and Continuation}

Households have CRRA period utility,
\begin{equation}
  u(c)=\frac{c^{1-\sigma}}{1-\sigma}.
  \whenintegrated{\label{eq:utility}}
\end{equation}
Future utility is discounted by $\beta$. The model also includes intergenerational altruism through the parameter $\eta$. When $\eta=0$, the household has no voluntary bequest motive. When $\eta=1$, the household values the descendant like a dynastic continuation.

Demographics are represented with two survival probabilities. A young household remains young with probability $\pi_y$ and becomes old with probability $1-\pi_y$. An old household survives with probability $\pi_o$ and dies with probability $1-\pi_o$, in which case a descendant enters the model with inherited assets. The notebook keeps these continuation weights in the Bellman equations, although its stationary distribution diagnostic is young-only rather than a full demographic distribution.

\subsection{Worker Income and Entrepreneurial Technology}

Workers earn labor income from their worker productivity:
\begin{equation}
  a' = (1+r)a + (1-\tau)wy - c.
  \whenintegrated{\label{eq:worker-budget}}
\end{equation}
Here $r$ is the risk-free interest rate, $w$ is the wage, and $\tau$ is the payroll tax used to finance pensions. A retired household receives pension income $p$ and solves a similar consumption-saving problem:
\begin{equation}
  a' = (1+r)a + p - c.
  \whenintegrated{\label{eq:retiree-budget}}
\end{equation}

Entrepreneurs operate a decreasing-returns technology. If an entrepreneur invests working capital $k$, output is
\[
  \theta k^\nu,
\]
with $\nu<1$. The entrepreneur's resources after production and debt repayment are
\begin{equation}
  a' = (1-\delta)k + \theta k^\nu - (1+r)(k-a) - c.
  \whenintegrated{\label{eq:budget-entrepreneur}}
\end{equation}
This equation is the main reason entrepreneurship is different from ordinary saving. The entrepreneur can choose a business scale $k$ that may exceed assets $a$, but doing so requires borrowing and therefore must satisfy the enforcement constraint.

\subsection{Borrowing and Enforcement}

The model does not impose a simple fixed borrowing limit. Instead, borrowing is limited by imperfect contract enforcement. If an entrepreneur defaults, the entrepreneur keeps a fraction $f$ of working capital and becomes a worker for one period. Lenders can seize the remaining fraction. Therefore a feasible entrepreneurial choice must make repayment at least as valuable as default.

For a young entrepreneur the enforcement condition is
\begin{equation}
  u(c)+\beta\pi_y E V(a',y',\theta')
  +\beta(1-\pi_y)E W(a',\theta')
  \geq V_w(fk,y,\theta).
  \whenintegrated{\label{eq:young-enforcement}}
\end{equation}
The left side is the value of honoring the contract. The right side is the value of defaulting, keeping $fk$, and entering the worker branch. A larger $f$ means default is more attractive, so the borrowing constraint becomes tighter. In the notebook this condition is implemented by checking each candidate $k$ against the relevant default value.

\subsection{Worker and Entrepreneur Branches}

The young household first compares the worker and entrepreneur branches:
\begin{equation}
  V(a,y,\theta)=\max\{V_w(a,y,\theta),V_e(a,y,\theta)\}.
  \whenintegrated{\label{eq:young-value}}
\end{equation}
The worker branch is a standard consumption-saving problem:
\begin{equation}
  V_w(a,y,\theta)
  = \max_{c,a'}
  \left\{
    u(c)+\beta\pi_y E V(a',y',\theta')
    +\beta(1-\pi_y) W_r(a')
  \right\},
  \whenintegrated{\label{eq:young-worker}}
\end{equation}
subject to equation~\eqref{eq:worker-budget} and $a'\geq 0$. If a young worker ages, the model sends the household to retirement, so the continuation term uses $W_r$.

The entrepreneur branch adds the business capital choice:
\begin{equation}
  V_e(a,y,\theta)
  = \max_{c,k,a'}
  \left\{
    u(c)+\beta\pi_y E V(a',y',\theta')
    +\beta(1-\pi_y)E W(a',\theta')
  \right\},
  \whenintegrated{\label{eq:young-entrepreneur}}
\end{equation}
subject to equation~\eqref{eq:budget-entrepreneur}, $a'\geq 0$, $k\geq 0$, and the enforcement condition~\eqref{eq:young-enforcement}. This is the branch where the reproduction has to do the most work, because $k$ affects both business returns and the default value.

\subsection{Old Households, Retirement, and Bequests}

Old entrepreneurs choose whether to continue operating the business or retire:
\begin{equation}
  W(a,\theta)=\max\{W_e(a,\theta),W_r(a)\}.
  \whenintegrated{\label{eq:old-value}}
\end{equation}
The continuing entrepreneur branch is similar to the young entrepreneur branch, but the continuation value reflects survival as old and possible death with a descendant:
\begin{equation}
  W_e(a,\theta)=\max_{c,k,a'}
  \left\{
    u(c)+\beta\pi_o E W(a',\theta')
    +\eta\beta(1-\pi_o)E V(a',y',\theta')
  \right\}.
  \whenintegrated{\label{eq:old-entrepreneur}}
\end{equation}
The old entrepreneur faces the same resource equation and a default comparison against retirement value $W_r(fk)$.

Retirement is absorbing in the sense that a retired old household does not re-enter entrepreneurship:
\begin{equation}
  W_r(a)=\max_{c,a'}
  \left\{
    u(c)+\beta\pi_o W_r(a')
    +\eta\beta(1-\pi_o)E V(a',y',\theta')
  \right\},
  \whenintegrated{\label{eq:retiree}}
\end{equation}
subject to equation~\eqref{eq:retiree-budget}. This equation is also where voluntary bequests enter. When old households die, descendants receive assets and draw ability states, so saving can matter for both the current household and the next generation.

\subsection{Stationary Equilibrium}
\whenintegrated{\label{sec:equilibrium}}

The full paper closes the model in stationary general equilibrium. A stationary equilibrium consists of prices $(r,w)$, a tax rate $\tau$, policy functions for consumption, saving, occupation, and entrepreneurial investment, and an invariant distribution over household states. Prices are consistent with the nonentrepreneurial production sector, the social-security budget balances, and the distribution is unchanged by the household transition rules.

For this reproduction, I only use that equilibrium definition to understand what would be needed for a full replication. The current notebook fixes prices and policy inputs, solves the household problem, and reports diagnostics from that partial-equilibrium solution. Therefore the theory reproduction reconstructs the main dynamic program, but it does not claim to reproduce the paper's full stationary general equilibrium.

\smartbib

\end{document}
